\chapter{Kesimpulan dan Saran}
    \section{Kesimpulan}

        Penelitian ini mengimplementasikan pendekatan \textit{two-stage} untuk prediksi pendapatan layanan SMS internasional pada perusahaan X menggunakan \textit{tree-based models}. Pada \textit{Stage} 1, model klasifikasi yaitu \textit{XGBoost} digunakan untuk memprediksi status aktif(pendapatan tidak nol) atau tidaknya (pendapatan nol) suatu \texttt{Series\_ID} pada bulan berikutnya. Sedangkan pada \textit{Stage} 2, model regresi yaitu \textit{Random Forest} dan \textit{XGBoost} dengan optimasi hiperparameter menggunakan Algoritma Genetika digunakan untuk memprediksi nilai pendapatan pada \texttt{Series\_ID} yang aktif saja berdasarkan hasil prediksi dari \textit{Stage} 1. 

        Penelitian ini menunjukkan bahwa optimasi hyperparameter menggunakan Algoritma Genetika dapat meningkatkan performansi dari \textit{tree-based models}. Metode ini terbukti efektif untuk meningkatkan prediksi \textit{tree-based Models} pada kasus prediksi pendapatan layanan SMS Internasional pada perusahaan X. Hasil menunjukkan bahwa XGBoost mengungguli model yang lain yaitu \textit{Random Forest} untuk performansi prediksi pendapatan layanan SMS Internasional pada perusahaan X dengan nilai WMAPE sebesar 16.29\%, MAE sebesar 23.398, 93 dan R\^2 sebesar 0.9765. Hasil ini memilik peningkatan siginifikan pada performansi model dengan penurunan WMAPE sebesar 10.83 \% dibandingkan dengan model baselinenya. Model XGBoost juga menampilkan tingkat konvergensi yang baik pada tiap generasinya. Akan tetapi, penelitian ini belum membandingkan hasil Algoritma Genetika dengan algoritma metaheuristik lain seperti \textit{Particle Swarm Optimization} (PSO), serta belum mengeksplorasi penerapan Algoritma Genetika untuk optimasi model \textit{deep learning}.


    \section{Saran}

        Setelah melakukan berbagai macam eksperimen pada penelitian ini, penulis menyarankan beberapa hal untuk penelitian selanjutnya.  Pertama, membandingkan dengan algoritma mataheuristik yang lain seperti PSO (\textit{Particle Swarm Optimization}) atau \textit{Bayesian Optimization} untuk mengetahui performansi algortima genetika dibanding dengan metodologi optimasi hiperparameter lain. Kedua, mengeksplorasi penambahan fitur eksternal seperti indikator ekonomi, pola musiman, atau kebijakan regulasi perusahaan telekomunikasi internasional untuk meningkatkan akurasi model. Terakhir, membandingkan pendekatan \textit{tree-based models} dengan model \textit{neural network} untuk mengetahui pendekatan yang lebih kompleks mampu memberikan hasil prediksi yang lebih baik. 
